\documentclass{article}
\usepackage[T1]{fontenc}
\usepackage{lmodern}
\usepackage{graphicx}
\usepackage{amsmath,amssymb}
\usepackage[a4paper,margin=2cm]{geometry}
\usepackage[absolute,overlay]{textpos}
\setlength{\TPHorizModule}{1mm}
\setlength{\TPVertModule}{1mm}
\usepackage[indonesian]{babel}
\usepackage{hyperref}
\setlength{\emergencystretch}{2em}
% unit: Halaman 45

\begin{document}

% === Halaman 45 (python/native) ===

45

3.

Lakukan pengacakan (permutasi) nilai-nilai di dalam larik S  

berdasarkan kunci rahasia K (kunci eksternal dari pengguna):


j  0


for i  0 to 255 do

       j  (j + S[i] + K[i  mod Length(K)]) mod 256

       swap(S[i], S[j]) \{* Pertukarkan S[i] \& S[j] *\}

     end

• Permutasi ini menyebabkan elemen-elemen di dalam larik S teracak. 

• K[i mod Length(K)] menyatakan karakter-karakter kunci diulang 

secara periodik jika panjangnya kurang dari 256

\end{document}
